\section{Całka Darboux}
% Import: Analiza matematyczna I, rozdział 29 i rozdział 19 (zachowane etykiety).
\begin{definicja}
\label{ma1:def:riemann-partition}
Niech $a<b$. Podziałem przedziału $[a,b]$ nazywamy skończony ciąg $P=(x_0,\ldots,x_n)$, gdzie $n\ge1$ oraz $a=x_0<x_1<\cdots<x_n=b$. Oznaczamy $\Delta x_i=x_i-x_{i-1}$ oraz
\[
 |P|=\max_{1\le i\le n}\Delta x_i.
\]
Liczbę $|P|$ (oznaczaną też przez $\mesh P$) nazywamy średnicą podziału. Podział $Q$ jest rozdrobnieniem $P$, jeśli zawiera wszystkie jego punkty. Suma zbiorów punktów dwóch podziałów, uporządkowana rosnąco, jest ich wspólnym rozdrobnieniem.
\end{definicja}
\begin{definicja}
\label{ma1:def:darboux-sums}
Niech $f:[a,b]\to\R$ będzie ograniczona. Dla $I_i=[x_{i-1},x_i]$ definiujemy
\[
 m_i=\inf_{I_i}f,\qquad M_i=\sup_{I_i}f,\qquad
 L(f,P)=\sum_{i=1}^n m_i\Delta x_i,\qquad
 U(f,P)=\sum_{i=1}^n M_i\Delta x_i.
\]
Są to odpowiednio suma dolna i suma górna Darboux. Ich różnica wynosi $\sum_i\omega_i\Delta x_i$, gdzie $\omega_i=M_i-m_i$ jest oscylacją $f$ na $I_i$.
\end{definicja}
Kresy istnieją dzięki ograniczoności $f$ i zupełności $\R$. Nie muszą być osiągane. Na przykład funkcja $f(x)=x$ dla $0<x<1$, z wartościami $f(0)=f(1)=1/2$, ma na $[0,1]$ kres dolny $0$ i górny $1$: wartości $f(1/n)$ i $f(1-1/n)$ dla $n\ge2$ dążą odpowiednio do $0$ i $1$, lecz wszystkie wartości funkcji należą do $(0,1)$. Jeżeli $f$ jest nieujemna, sumy można interpretować jako pola prostokątów zawartych pod wykresem i pokrywających obszar pod wykresem. Dla funkcji zmieniających znak są to sumy ze znakiem.

\begin{center}
\begin{tikzpicture}[x=5.2cm,y=3.2cm,>=Stealth]
 \draw[->] (-.07,0)--(1.12,0) node[right] {$x$};
 \draw[->] (0,-.07)--(0,1.15) node[above] {$y$};
 \foreach \i in {0,1,2,3}{
   \pgfmathsetmacro{\l}{\i/4}
   \pgfmathsetmacro{\r}{(\i+1)/4}
   \pgfmathsetmacro{\lo}{\l*\l}
   \pgfmathsetmacro{\hi}{\r*\r}
   \fill[blue!18] (\l,0) rectangle (\r,\lo);
   \fill[orange!30] (\l,\lo) rectangle (\r,\hi);
   \draw[orange!80!black] (\l,0) rectangle (\r,\hi);
   \draw[blue!65!black] (\l,\lo)--(\r,\lo);
 }
 \draw[very thick,black,domain=0:1,samples=100] plot (\x,{\x*\x});
 \foreach \i in {1,2,3,4}{
   \pgfmathsetmacro{\t}{\i/4}
   \draw (\t,-.018)--(\t,.018);
   \node[below] at (\t,-.02) {$x_{\i}$};
 }
 \node[below left] at (0,0) {$x_0$};
 \node[above left] at (.88,.83) {$f(x)=x^2$};
\end{tikzpicture}

\small Rysunek: suma dolna (niebieska) i suma górna (niebieska wraz z pomarańczową częścią) dla $f(x)=x^2$ i czterech równych podprzedziałów $[0,1]$. Pomarańczowa część przedstawia różnicę sum.
\end{center}

\begin{lemat}
\label{ma1:lem:darboux-refinement}
Jeżeli $Q$ jest rozdrobnieniem $P$, to
\[
 L(f,P)\le L(f,Q)\le U(f,Q)\le U(f,P).
\]
Ponadto dla dowolnych podziałów $P,R$ zachodzi $L(f,P)\le U(f,R)$.
\end{lemat}
\begin{proof}
Niech $f$ będzie ograniczona na $[a,b]$. Każdy podprzedział $P$ jest sumą kolejnych podprzedziałów $Q$. Kres dolny na każdym mniejszym przedziale jest nie mniejszy niż kres dolny na większym, a kres górny jest nie większy niż kres górny na większym. Mnożymy te nierówności przez dodatnie długości mniejszych przedziałów i sumujemy; suma ich długości jest długością większego przedziału. Otrzymujemy zewnętrzne nierówności. Środkowa wynika z $m_i\le M_i$. Dla dowolnych $P,R$ bierzemy wspólne rozdrobnienie $Q$; wtedy $L(f,P)\le L(f,Q)\le U(f,Q)\le U(f,R)$.
\end{proof}
\begin{definicja}
\label{ma1:def:riemann-integral}
Dla ograniczonej funkcji $f:[a,b]\to\R$ definiujemy całkę dolną i górną, biorąc kresy po wszystkich podziałach $[a,b]$:
\[
 \underline{\mathcal I}(f)=\sup_P L(f,P),\qquad
 \overline{\mathcal I}(f)=\inf_P U(f,P).
\]
Jeżeli te liczby są równe, mówimy, że $f$ jest całkowalna w sensie Darboux na $[a,b]$, a wspólną wartość oznaczamy przez $\int_a^b f(x)\,\dd x$.
\end{definicja}
Jeśli $m\le f\le M$, to wszystkie sumy leżą w $[m(b-a),M(b-a)]$. Oba kresy są więc skończone. Z lematu \ref{ma1:lem:darboux-refinement}, najpierw biorąc supremum po $P$, a następnie infimum po $R$, dostajemy $\underline{\mathcal I}(f)\le\overline{\mathcal I}(f)$.

\begin{twierdzenie}[Kryterium Darboux]
\label{ma1:thm:darboux-criterion}
Funkcja ograniczona $f$ jest całkowalna w sensie Darboux na $[a,b]$ wtedy i tylko wtedy, gdy dla każdego $\varepsilon>0$ istnieje podział $P$ taki, że $U(f,P)-L(f,P)<\varepsilon$.
\end{twierdzenie}
\begin{proof}
Jeżeli obie całki są równe $A$, z definicji kresów wybieramy $P,R$ takie, że $L(f,P)>A-\varepsilon/2$ i $U(f,R)<A+\varepsilon/2$. Dla wspólnego rozdrobnienia $Q$ mamy $U(f,Q)-L(f,Q)\le U(f,R)-L(f,P)<\varepsilon$. Odwrotnie, dla każdego podziału $P$ mamy
\[
 0\le\overline{\mathcal I}(f)-\underline{\mathcal I}(f)
 \le U(f,P)-L(f,P).
\]
Jeżeli prawa strona może być mniejsza od dowolnego $\varepsilon>0$, różnica całek jest zerowa.
\end{proof}

\subsection{Kryteria całkowalności i przykłady}
\begin{twierdzenie}
\label{ma1:thm:riemann-continuous}
Każda funkcja ciągła $f:[a,b]\to\R$ jest całkowalna.
\end{twierdzenie}
\begin{proof}
Funkcja jest ograniczona z twierdzenia Weierstrassa i jednostajnie ciągła z twierdzenia Heinego--Cantora. Dla $\varepsilon>0$ istnieje $\delta>0$ takie, że $|f(s)-f(t)|<\varepsilon/(2(b-a))$ dla $|s-t|<\delta$. Dla podziału o $|P|<\delta$ oscylacja na każdym $I_i$ jest nie większa niż $\varepsilon/(2(b-a))$ (możemy wziąć punkty osiągania maksimum i minimum). Stąd $U(f,P)-L(f,P)\le\varepsilon/2<\varepsilon$. Stosujemy kryterium Darboux.
\end{proof}
\begin{twierdzenie}
\label{ma1:thm:riemann-monotone}
Każda funkcja monotoniczna $f:[a,b]\to\R$ jest całkowalna.
\end{twierdzenie}
\begin{proof}
Jeśli $f$ jest niemalejąca, to $m_i=f(x_{i-1})$, $M_i=f(x_i)$, a wartości funkcji leżą pomiędzy $f(a)$ i $f(b)$. Zatem
\[
 U(f,P)-L(f,P)\le |P|\sum_i(f(x_i)-f(x_{i-1}))
 =|P|(f(b)-f(a)).
\]
Podziały o dowolnie małej średnicy dają kryterium Darboux. Dla funkcji nierosnącej mamy $m_i=f(x_i)$ i $M_i=f(x_{i-1})$; ten sam rachunek daje oszacowanie $|P|(f(a)-f(b))$.
\end{proof}
\begin{twierdzenie}
\label{ma1:thm:riemann-finite-discontinuities}
Funkcja ograniczona na $[a,b]$ i ciągła poza skończonym zbiorem punktów jest całkowalna.
\end{twierdzenie}
\begin{proof}
Niech $|f|\le B$ dla $B>0$, a zbiór możliwych nieciągłości ma $k$ punktów. Gdy $k=0$, stosujemy twierdzenie \ref{ma1:thm:riemann-continuous}. Gdy $k>0$, oznaczamy te punkty przez $d_1,\ldots,d_k$. Wybieramy $\rho>0$ mniejsze od $(b-a)/2$ oraz $\varepsilon/(8Bk)$, a gdy $k\ge2$, także od połowy najmniejszej odległości między różnymi $d_j$. Dla każdego $j$ bierzemy
\[
 J_j=[\max(a,d_j-\rho),\min(b,d_j+\rho)].
\]
Każdy $J_j$ ma dodatnią długość, zawiera $d_j$ we wnętrzu względem $[a,b]$ i nie przecina żadnego innego $J_i$. Suma długości jest nie większa niż $2k\rho<\varepsilon/(4B)$. Włączamy do podziału $a,b$ oraz wszystkie końce $J_j$. Na nich łączny wkład do różnicy sum jest mniejszy od $2B\varepsilon/(4B)=\varepsilon/2$.

Pozostałe odcinki o dodatniej długości pomiędzy tymi przedziałami są domknięte i nie zawierają punktów nieciągłości: końce wewnętrzne wybranych przedziałów nie są tymi punktami. Na każdym z tych odcinków funkcja jest ciągła. Jest ich skończenie wiele, powiedzmy $\ell$. Jeżeli $\ell>0$, wybieramy na każdym podział z różnicą sum mniejszą od $\varepsilon/(2\ell)$ na mocy twierdzenia \ref{ma1:thm:riemann-continuous} i kryterium Darboux. Łączymy wszystkie podziały. Wkłady na pozostałych odcinkach sumują się do mniej niż $\varepsilon/2$; gdy $\ell=0$, ich wkład wynosi zero. Cała różnica sum jest mniejsza od $\varepsilon$.
\end{proof}
\begin{przyklad}
\label{ma1:ex:riemann-dirichlet}
Funkcja $\chi_{\Q\cap[a,b]}$ nie jest całkowalna, gdy $a<b$. Każdy podprzedział dodatniej długości zawiera liczbę wymierną oraz niewymierną, więc jego kresy wynoszą $0$ i $1$. Dla przypomnienia: między $u<v$ wybieramy liczbę całkowitą $n$ z $n(v-u)>1$, a następnie całkowite $m$ z $nu<m<nv$; wtedy $m/n\in(u,v)$. Wybierając wymierną w $(u-\sqrt2,v-\sqrt2)$ i dodając $\sqrt2$, otrzymujemy liczbę niewymierną w $(u,v)$. Stąd dla każdego podziału $L=0$, $U=b-a$, co przeczy kryterium Darboux. 
\end{przyklad}

\begin{przyklad}[Funkcja Thomae'a]
\label{ma1:ex:thomae}
Definiujemy $T:\R\to[0,1]$ przez
\[
 T(x)=
 \begin{cases}
 1/q,& x=p/q,\ p\in\Z,\ q\in\N,\ q\ge1,\ \gcd(|p|,q)=1,\\
 0,&x\notin\Q.
 \end{cases}
\]
Ułamek nieskracalny o dodatnim mianowniku jest jednoznaczny; w szczególności $T(0)=1$. Dla każdego $a\in\R$ istnieje granica przekłuta $\lim_{x\to a}T(x)=0$. Zatem $T$ jest ciągła dokładnie w punktach niewymiernych.
\end{przyklad}
\begin{proof}
Ustalmy $a\in\R$ i $\varepsilon>0$. Wybierzmy całkowite $N\ge1$ z $1/N<\varepsilon$. Zbiór
$E_N=\{p/q\in[a-1,a+1]:p\in\Z,\ 1\le q\le N,\ \gcd(|p|,q)=1\}$
jest skończony: dla każdego spośród skończenie wielu mianowników $q$ licznik należy do ograniczonego przedziału $[q(a-1),q(a+1)]$. Przyjmijmy $\delta$ jako minimum liczby $1$ i wszystkich dodatnich odległości $|r-a|$ dla $r\in E_N\setminus\{a\}$. Jeśli ten ostatni zbiór jest pusty, przyjmujemy $\delta=1$. Mamy $\delta>0$.

Dla $0<|x-a|<\delta$ punkt $x$ należy do $[a-1,a+1]$, ale nie do $E_N$. Jeśli $x$ jest niewymierny, to $T(x)=0$; jeśli jest wymierny, jego nieskracalny mianownik jest większy od $N$, więc $0<T(x)<1/N<\varepsilon$. Dowodzi to granicy przekłutej. Ponieważ każdy punkt $\R$ jest punktem skupienia dziedziny, ciągłość jest równoważna $T(a)=0$, czyli $a\notin\Q$.
\end{proof}
\begin{center}
\begin{tikzpicture}[x=9cm,y=2.7cm,>=Stealth]
 \draw[->] (-.03,0)--(1.08,0) node[right] {$x$};
 \draw[->] (0,-.08)--(0,1.15) node[above] {$y$};
 \fill[blue!65!black] ({0/1},{1/1}) circle (.9pt);
 \fill[blue!65!black] ({1/1},{1/1}) circle (.9pt);
 \fill[blue!65!black] ({1/2},{1/2}) circle (.9pt);
 \fill[blue!65!black] ({1/3},{1/3}) circle (.9pt);
 \fill[blue!65!black] ({2/3},{1/3}) circle (.9pt);
 \fill[blue!65!black] ({1/4},{1/4}) circle (.9pt);
 \fill[blue!65!black] ({3/4},{1/4}) circle (.9pt);
 \fill[blue!65!black] ({1/5},{1/5}) circle (.9pt);
 \fill[blue!65!black] ({2/5},{1/5}) circle (.9pt);
 \fill[blue!65!black] ({3/5},{1/5}) circle (.9pt);
 \fill[blue!65!black] ({4/5},{1/5}) circle (.9pt);
 \fill[blue!65!black] ({1/6},{1/6}) circle (.9pt);
 \fill[blue!65!black] ({5/6},{1/6}) circle (.9pt);
 \fill[blue!65!black] ({1/7},{1/7}) circle (.9pt);
 \fill[blue!65!black] ({2/7},{1/7}) circle (.9pt);
 \fill[blue!65!black] ({3/7},{1/7}) circle (.9pt);
 \fill[blue!65!black] ({4/7},{1/7}) circle (.9pt);
 \fill[blue!65!black] ({5/7},{1/7}) circle (.9pt);
 \fill[blue!65!black] ({6/7},{1/7}) circle (.9pt);
 \fill[blue!65!black] ({1/8},{1/8}) circle (.9pt);
 \fill[blue!65!black] ({3/8},{1/8}) circle (.9pt);
 \fill[blue!65!black] ({5/8},{1/8}) circle (.9pt);
 \fill[blue!65!black] ({7/8},{1/8}) circle (.9pt);
 \fill[blue!65!black] ({1/9},{1/9}) circle (.9pt);
 \fill[blue!65!black] ({2/9},{1/9}) circle (.9pt);
 \fill[blue!65!black] ({4/9},{1/9}) circle (.9pt);
 \fill[blue!65!black] ({5/9},{1/9}) circle (.9pt);
 \fill[blue!65!black] ({7/9},{1/9}) circle (.9pt);
 \fill[blue!65!black] ({8/9},{1/9}) circle (.9pt);
 \fill[blue!65!black] ({1/10},{1/10}) circle (.9pt);
 \fill[blue!65!black] ({3/10},{1/10}) circle (.9pt);
 \fill[blue!65!black] ({7/10},{1/10}) circle (.9pt);
 \fill[blue!65!black] ({9/10},{1/10}) circle (.9pt);
 \fill[blue!65!black] ({1/11},{1/11}) circle (.9pt);
 \fill[blue!65!black] ({2/11},{1/11}) circle (.9pt);
 \fill[blue!65!black] ({3/11},{1/11}) circle (.9pt);
 \fill[blue!65!black] ({4/11},{1/11}) circle (.9pt);
 \fill[blue!65!black] ({5/11},{1/11}) circle (.9pt);
 \fill[blue!65!black] ({6/11},{1/11}) circle (.9pt);
 \fill[blue!65!black] ({7/11},{1/11}) circle (.9pt);
 \fill[blue!65!black] ({8/11},{1/11}) circle (.9pt);
 \fill[blue!65!black] ({9/11},{1/11}) circle (.9pt);
 \fill[blue!65!black] ({10/11},{1/11}) circle (.9pt);
 \node[below] at (0,0) {$0$};
 \node[below] at (1,0) {$1$};
 \node[left] at (0,1) {$1$};
 \node[blue!65!black,above] at (.5,.5) {$1/2$};
\end{tikzpicture}

\small Rysunek: wartości funkcji Thomae'a dla nieskracalnych mianowników nie większych niż $11$; pominięto wartości o większych mianownikach i wartości zerowe w punktach niewymiernych.
\end{center}

\begin{przyklad}[Całka Darboux funkcji Thomae'a]
Rozważamy ograniczenie funkcji z przykładu \ref{ma1:ex:thomae} do $[0,1]$.
Dla dowolnego podziału $P$ mamy $L(T,P)=0$, gdyż każdy podprzedział zawiera liczbę niewymierną, a $T\ge0$. Ustalmy $\varepsilon>0$ i wybierzmy $N$ z $1/N<\varepsilon/2$. Niech $E_N$ będzie zbiorem liczb wymiernych w $[0,1]$ o mianowniku nie większym niż $N$ w zapisie nieskracalnym; jest skończony z poprzedniego dowodu. Oznaczmy $k=|E_N|$; mamy $k\ge1$. Wybieramy podział $P$ o średnicy mniejszej niż $\varepsilon/(4k)$. Każdy punkt $E_N$ należy do co najwyżej dwóch domkniętych podprzedziałów, dlatego podprzedziałów przecinających $E_N$ jest co najwyżej $2k$, a suma ich długości nie przekracza $2k\mesh P$. Na nich $T\le1$, zaś na pozostałych $T\le1/N$. Otrzymujemy
\[
0\le U(T,P)-L(T,P)=U(T,P)
\le 2k\mesh P+\frac1N<\varepsilon.
\]
Z kryterium Darboux wynika całkowalność w sensie Darboux. Ponieważ każda suma dolna wynosi zero, wspólna wartość obu całek Darboux jest równa zeru.
\end{przyklad}
\begin{przyklad}[Obliczenie z definicji]
Dla $f(x)=x^2$ na $[0,1]$ funkcja jest niemalejąca, zatem na $[(i-1)/n,i/n]$ infimum wynosi $((i-1)/n)^2$, a supremum $(i/n)^2$. Korzystamy ze wzoru $\sum_{i=1}^n i^2=n(n+1)(2n+1)/6$: prawa strona dla $n=0$ wynosi zero, a jej przyrost od $n-1$ do $n$ to
\[
\frac{n}{6}\bigl((n+1)(2n+1)-(n-1)(2n-1)\bigr)=n^2,
\]
co dowodzi wzoru przez sumowanie przyrostów. Dla $P_n=\{i/n:0\le i\le n\}$ mamy
\[
U(f,P_n)=\frac1{n^3}\sum_{i=1}^n i^2=\frac{(n+1)(2n+1)}{6n^2},\qquad
L(f,P_n)=\frac1{n^3}\sum_{i=0}^{n-1}i^2=\frac{(n-1)(2n-1)}{6n^2}.
\]
Dzieląc liczniki przez $n^2$, otrzymujemy, że obie sumy dążą do $1/3$. Dla każdego $n$ suma dolna nie przekracza całki dolnej, ta nie przekracza całki górnej, a całka górna nie przekracza sumy górnej. Z twierdzenia o trzech ciągach (lub bezpośrednio z tych nierówności) obie całki Darboux wynoszą $1/3$, czyli $\int_0^1x^2\,\mathrm dx=1/3$.
\end{przyklad}
